% =============================================================================
% Statements of mock midterm 3 (Topics 1, 2A, 2B), written once and \input by
%   mock3-midterm-topics-1-2b.tex            (exam paper)
%   mock3-midterm-topics-1-2b-solutions.tex  (worked solutions)
% Every statement is a real exercise copied verbatim (rules/exercise-lists.md):
% source line first, then the full text. Allowed edits only: math in LaTeX, tables
% rebuilt, line breaks and hyphenation undone, currency written as \$.
% Needs: \srcline{...} defined by the including document.
% =============================================================================

% ---------------------------------------------------------------- Question 1.1
\newcommand{\QoneA}{%
\srcline{Higgins (Koski--Higgins--Mitton), \emph{Analysis for Financial Management} 13e, Ch.~2,
Exercise 3, pp.~71--72 (PDF 92--93) $\cdot$ Answer in book: yes, p.~401 (PDF 422)}

\noindent Answer true or false to each of the following. Briefly explain your reasoning for each answer.
\begin{enumerate}[label=\alph*.,leftmargin=2em]
  \item A company's assets-to-equity ratio always equals one plus its liabilities-to-equity ratio.
  \item A company's return on equity will always equal or exceed its return on assets.
  \item A company's collection period should always be less than its payables period.
  \item A company's current ratio must always equal or exceed its acid-test ratio.
  \item All else equal, a firm would prefer to have a higher asset turnover ratio.
  \item Two firms can have the same earnings yield but different price-to-earnings ratios.
  \item Ignoring taxes and transactions costs, unrealized paper gains are less valuable than realized
        cash earnings.
\end{enumerate}}

% ---------------------------------------------------------------- Question 1.2
\newcommand{\QoneB}{%
\srcline{Higgins (Koski--Higgins--Mitton), \emph{Analysis for Financial Management} 13e, Ch.~4,
Exercise 3, p.~134 (PDF 155) $\cdot$ Answer in book: yes, pp.~408--409 (PDF 429--430)}

\noindent Are the following statements true or false? Please explain why.
\begin{enumerate}[label=\alph*.,leftmargin=2em]
  \item The only way a company can grow at a rate above its current sustainable growth rate is by
        issuing new stock.
  \item The stock market is a ready source of new capital when a company is incurring heavy losses.
  \item Share repurchases usually increase earnings per share.
  \item Companies often buy back their stock because managers believe the shares are undervalued.
  \item Only rapidly growing firms have growth management problems.
  \item Increasing growth increases stock price.
\end{enumerate}}

% ---------------------------------------------------------------- Question 1.3
\newcommand{\QoneC}{%
\srcline{Berk, DeMarzo \& Harford, \emph{Fundamentals of Corporate Finance} 2e, Ch.~16, Critical
Thinking question 2, p.~490 (PDF 521) $\cdot$ Answer in book: no}

\noindent Explain what is wrong with the following argument: ``If a firm issues debt that is risk-free
because there is no possibility of default, the risk of the firm's equity does not change. Therefore,
risk-free debt allows the firm to get the benefit of a low cost of capital of debt without raising its
cost of capital of equity.''}

% ---------------------------------------------------------------- Question 2 (section instruction)
\newcommand{\QtwoIntro}{%
\noindent\emph{Section instruction} (Berk, Ch.~16, ``Capital Structure in Perfect Capital Markets'',
p.~490, PDF 521, applying to Parts 2.1, 2.2 and 2.4): For problems in this section assume no taxes or
distress costs.}

% ---------------------------------------------------------------- Question 2.1
\newcommand{\QtwoA}{%
\srcline{Berk, DeMarzo \& Harford, \emph{Fundamentals of Corporate Finance} 2e, Ch.~16, Exercise 16.2,
p.~491 (PDF 522) $\cdot$ Answer in book: no}

\noindent You are an entrepreneur starting a biotechnology firm. If your research is successful, the
technology can be sold for \$30 million. If your research is unsuccessful, it will be worth nothing. To
fund your research, you need to raise \$2 million. Investors are willing to provide you with \$2 million
in initial capital in exchange for 50\% of the unlevered equity in the firm.
\begin{enumerate}[label=\alph*.,leftmargin=2em]
  \item What is the total market value of the firm without leverage?
  \item Suppose you borrow \$1 million. According to MM, what fraction of the firm's equity will you
        need to sell to raise the additional \$1 million you need?
  \item What is the value of your share of the firm's equity in cases (a) and (b)?
\end{enumerate}}

% ---------------------------------------------------------------- Question 2.2
\newcommand{\QtwoB}{%
\srcline{Berk, DeMarzo \& Harford, \emph{Fundamentals of Corporate Finance} 2e, Ch.~16, Exercise *16.4,
p.~491 (PDF 522) $\cdot$ Answer in book: no}

\noindent Suppose there are no taxes. Firm ABC has no debt, and firm XYZ has debt of \$5000 on which it
pays interest of 10\% each year. Both companies have identical projects that generate free cash flows of
\$800 or \$1000 each year. After paying any interest on debt, both companies use all remaining free cash
flows to pay dividends each year.
\begin{center}
\begin{tabular}{lcccc}
  \toprule
  & \multicolumn{2}{c}{ABC} & \multicolumn{2}{c}{XYZ} \\
  \cmidrule(lr){2-3}\cmidrule(lr){4-5}
  FCF & Debt Payments & Equity Dividends & Debt Payments & Equity Dividends \\ \midrule
  \$\ 800 & & & & \\
  \$1000 & & & & \\ \bottomrule
\end{tabular}
\end{center}
\begin{enumerate}[label=\alph*.,leftmargin=2em]
  \item Fill in the table above showing the payments debt and equity holders of each firm will receive
        given each of the two possible levels of free cash flows.
  \item Suppose you hold 10\% of the equity of ABC. What is another portfolio you could hold that would
        provide the same cash flows?
  \item Suppose you hold 10\% of the equity of XYZ. If you can borrow at 10\%, what is an alternative
        strategy that would provide the same cash flows?
\end{enumerate}}

% ---------------------------------------------------------------- Question 2.3
\newcommand{\QtwoC}{%
\srcline{Berk, DeMarzo \& Harford, \emph{Fundamentals of Corporate Finance} 2e, Ch.~13, Exercise 13.21,
p.~405 (PDF 436) $\cdot$ Answer in book: no}

\noindent Your company has two divisions: One division sells software and the other division sells
computers through a direct sales channel, primarily taking orders over the Internet. You have decided
that Dell Computer is very similar to your computer division, in terms of both risk and financing. You
go online and find the following information: Dell's beta is 1.21, the risk-free rate is 4.5\%, its
market value of equity is \$67 billion, and it has \$700 million worth of debt with a yield to maturity
of 6\%. Your tax rate is 35\% and you use a market risk premium of 5\% in your WACC estimates.
\begin{enumerate}[label=\alph*.,leftmargin=2em]
  \item What is an estimate of the WACC for your computer sales division?
  \item If your overall company WACC is 12\% and the computer sales division represents 40\% of the
        value of your firm, what is an estimate of the WACC for your software division?
\end{enumerate}}

% ---------------------------------------------------------------- Question 2.4
\newcommand{\QtwoD}{%
\srcline{Berk, DeMarzo \& Harford, \emph{Fundamentals of Corporate Finance} 2e, Ch.~16, Exercise 16.6,
p.~492 (PDF 523) $\cdot$ Answer in book: no}

\noindent Suppose Microsoft has no debt and a WACC of 9.2\%. The average debt-to-value ratio for the
software industry is 5\%. What would its cost of equity be if it took on the average amount of debt for
its industry at a cost of debt of 6\%?}

% ---------------------------------------------------------------- Question 3.1
\newcommand{\QthreeA}{%
\srcline{Berk, DeMarzo \& Harford, \emph{Fundamentals of Corporate Finance} 2e, Ch.~16, Exercises
16.7--16.8, p.~492 (PDF 523) $\cdot$ Answer in book: no}

\noindent\textbf{7.} Pelamed Pharmaceuticals had EBIT of \$325 million in 2010. In addition, Pelamed had
interest expenses of \$125 million and a corporate tax rate of 40\%.
\begin{enumerate}[label=\alph*.,leftmargin=2em]
  \item What was Pelamed's 2010 net income?
  \item What was the total of Pelamed's 2010 net income and interest payments?
  \item If Pelamed had no interest expenses, what would its 2010 net income have been? How does it
        compare to your answer in part (b)?
  \item What was the amount of Pelamed's interest tax shield in 2010?
\end{enumerate}
\noindent\textbf{8.} Grommit Engineering expects to have net income next year of \$20.75 million and free
cash flow of \$22.15 million. Grommit's marginal corporate tax rate is 35\%.
\begin{enumerate}[label=\alph*.,leftmargin=2em]
  \item If Grommit increases leverage so that its interest expense rises by \$1 million, how will its
        net income change?
  \item For the same increase in interest expense, how will its free cash flow change?
\end{enumerate}}

% ---------------------------------------------------------------- Question 3.2
\newcommand{\QthreeB}{%
\srcline{Berk, DeMarzo \& Harford, \emph{Fundamentals of Corporate Finance} 2e, Ch.~16, Exercises 16.9
and 16.12, p.~492 (PDF 523) $\cdot$ Answer in book: no}

\noindent\textbf{9.} Assume that Microsoft has a total market value of \$300 billion and a marginal tax
rate of 35\%. If it permanently changes its leverage from no debt by taking on new debt in the amount of
13\% of its current market value, what is the present value of the tax shield it will create?

\smallskip\noindent\textbf{12.} Arnell Industries has \$10 million in permanent debt outstanding. The
firm will pay interest only on this debt. Arnell's marginal tax rate is expected to be 35\% for the
foreseeable future.
\begin{enumerate}[label=\alph*.,leftmargin=2em]
  \item Suppose Arnell pays interest of 6\% per year on its debt. What is its annual interest tax
        shield?
  \item What is the present value of the interest tax shield, assuming its risk is the same as the loan?
  \item Suppose instead that the interest rate on the debt is 5\%. What is the present value of the
        interest tax shield in this case?
\end{enumerate}}

% ---------------------------------------------------------------- Question 3.3
\newcommand{\QthreeC}{%
\srcline{Berk, DeMarzo \& Harford, \emph{Fundamentals of Corporate Finance} 2e, Ch.~16, Exercises
16.13--16.15, p.~493 (PDF 524) $\cdot$ Answer in book: no}

\noindent\textbf{13.} Rogot Instruments makes fine violins, violas, and cellos. It has \$1 million in debt
outstanding, equity valued at \$2 million, and pays corporate income tax at a rate of 35\%. Its cost of
equity is 12\% and its cost of debt is 7\%.
\begin{enumerate}[label=\alph*.,leftmargin=2em]
  \item What is Rogot's pretax WACC?
  \item What is Rogot's (effective after-tax) WACC?
\end{enumerate}
\noindent\textbf{14.} Rumolt Motors has 30 million shares outstanding with a price of \$15 per share. In
addition, Rumolt has issued bonds with a total current market value of \$150 million. Suppose Rumolt's
equity cost of capital is 10\%, and its debt cost of capital is 5\%.
\begin{enumerate}[label=\alph*.,leftmargin=2em]
  \item What is Rumolt's pretax WACC?
  \item If Rumolt's corporate tax rate is 35\%, what is its after-tax WACC?
\end{enumerate}
\noindent\textbf{15.} Summit Builders has a market debt-equity ratio of 0.65 and a corporate tax rate of
40\%, and it pays 7\% interest on its debt. By what amount does the interest tax shield from its debt
lower Summit's WACC?}

% ---------------------------------------------------------------- Question 3.4
\newcommand{\QthreeD}{%
\srcline{Berk, DeMarzo \& Harford, \emph{Fundamentals of Corporate Finance} 2e, Ch.~17, Exercises
17.3--17.5, p.~527 (PDF 558) $\cdot$ Answer in book: no. Problem 5 refers to Problem 4, so both are in
this block.}

\noindent\textbf{3.} ECB Co. has 1 million shares outstanding selling at \$20 per share. It plans to
repurchase 100,000 shares at the market price. What will its market capitalization be after the
repurchase? What will its stock price be?

\smallskip\noindent\textbf{4.} KMS corporation has assets of \$500 million, \$50 million of which are
cash. It has debt of \$200 million. If KMS repurchases \$20 million of its stock:
\begin{enumerate}[label=\alph*.,leftmargin=2em]
  \item What changes will occur on its balance sheet?
  \item What will its new leverage ratio be?
\end{enumerate}
\noindent\textbf{5.} Suppose that KMS in Problem 4 decides to initiate a dividend instead, but it wants
the present value of the payout to be the same \$20 million. If its cost of equity capital is 10\%, to
what amount per year in perpetuity should it commit (assuming perfect capital markets)?}
